% profiles/rutgers-university-new-brunswick.tex — Rutgers University-New Brunswick
% ============================================================================
% source: https://grad.rutgers.edu/academics/graduation/electronic-thesis-and-dissertation-style-guide
%         (Rutgers School of Graduate Studies, "Electronic Thesis and
%         Dissertation Style Guide")
% source: https://grad.rutgers.edu/sites/default/files/2024-02/title%20page.pdf
%         (the SGS "Doctorate Title Page" sample, linked from that guide)
% source: https://grad.rutgers.edu/sites/default/files/2024-02/samplemasterstitlepage.pdf
%         (the SGS "Master's Title Page" sample, linked from that guide)
% accessed: 2026-09-10
% checked:  texlib-thesis-profile-round
%
% VERIFIED and set here: geometry, spacing, title page, the absence of a
% separate approval page, and the accessibility requirement.
% NOT set: nothing is left to the neutral fallback.
%
% Rutgers is one graduate school (the School of Graduate Studies) serving
% several campuses, and the title page names the campus. This profile is the
% NEW BRUNSWICK one and prints "New Brunswick, New Jersey"; a Newark or Camden
% filing needs that line changed. (Rutgers-Camden's Graduate School publishes a
% separate style guide of its own, which was not read for this file.)

\thesisinstitution{Rutgers, The State University of New Jersey}

% "All margins must be 1 inch." Stated once, for all four sides, with no
% binding exception -- the manuscript is filed as a PDF through ProQuest.
\thesissetgeometry{left=1in, right=1in, top=1in, bottom=1in}

% "Double-space preliminary pages, appendices, and all text." The single-spaced
% items SGS lists -- footnotes, endnotes, references and bibliographic material,
% block quotations of four or more lines, verse quotations of two or more lines
% -- are per-element choices, not a document spacing rule.
\thesissetspacing{double}

% No chapter-opening exception: SGS states one margin figure and no deeper
% opening for chapters or major sections. \thesissetchapteropening is
% deliberately not called.

% Typeface: "Choose an easy-to-read type. Use one typeface throughout; script or
% italic typefaces are not acceptable for the main text (10-12 points)." No face
% is named, so this profile sets none; the class's Latin Modern satisfies it.

% --- Accessibility -----------------------------------------------------------
% SGS publishes a real requirement here, but it is FUTURE-DATED, and the modal
% verbs matter. Today the guide "encourages"; from 2027-04-26 it requires.
%
% NOTE FOR ANYONE COMPARING PROFILES: the date below is 26 April 2027. Other
% profiles in this directory carry 26 April 2026. Both are correct -- the ADA
% Title II compliance deadline differs by the size of the public entity -- and
% neither should be "corrected" against the other.
%
% \thesisrequiretagging is NOT set, and that is a judgement, not an oversight.
% SGS names Title II and points at its own PDF guide; it never says "tagged
% PDF", "PDF/UA" or "PDF/A" in the style guide. Declaring a tagging requirement
% here would be inferring the mechanism from the standard, which is exactly the
% inference RESEARCHING.md forbids. The requirement is recorded verbatim
% instead.
\thesisaccessibilityrequirement{Title II ADA, from 26 April 2027}
	{"Starting on April 26, 2027, all electronic dissertations and theses will
	 be required to comply with Title II ADA accessibility standards." Until
	 then the School of Graduate Studies "encourages authors to adhere to best
	 practices for accessibility in all newly submitted materials, including
	 electronic theses and dissertations."}
\thesisaccessibilityrequirement{Standard}
	{"The requirements for digital accessibility set by the School of Graduate
	 Studies are based on Title II of the ADA Accessibility Regulations." SGS
	 names no PDF standard of its own and refers authors to the SGS Quick Guide
	 on Accessibility Regulations and the SGS PDF Accessibility Guide.}
\thesisaccessibilityrequirement{Checking}
	{"You can verify compliance using Adobe Acrobat, but these checks can also
	 be performed in Microsoft Word before saving your document as a PDF."
	 SGS prescribes no check a LaTeX author can run; veraPDF is this class's
	 answer and is not Rutgers's.}

% --- Title page --------------------------------------------------------------
% Reproduced from the two SGS sample pages, which differ only in the words the
% class already varies by \doctype: "A dissertation submitted to the" against
% "A thesis submitted to the", and the degree line.
%
% The samples' own notes, verbatim: "Center and double space everything";
% "Title and author's name in upper case letters"; "Use your full, legal name";
% "Fill in the correct, official name of the program under 'Graduate program
% in'"; "Indicate your Dissertation Director's name [do not include their title]
% under 'Written under the direction of'"; "Number of lines corresponds to
% number of committee members"; "Indicate the month and the year the degree is
% to be conferred - October, January or May, 20XX-not the defense month and
% year."
%
% Two consequences for a document using this profile:
%   * \graddate must be one of October, January or May with the year of
%     CONFERRAL -- not the defense date. The profile cannot check that.
%   * \gradadvisor must carry the name alone. The class's default value ends in
%     ", Ph.D."; SGS says "do not use any formal title (Ph.D., Dr., etc.)".
%
% DOUBLE spacing on this page, not single: the class's other profiles switch to
% \singlespacing here, and SGS explicitly does not.
%
% The signature lines are BLANK BY REQUIREMENT, which is why the committee list
% is redefined to print a rule and no name: "The title page within your entire
% dissertation which you will upload to ProQuest for publishing must not be
% signed. That version has blank lines." One rule per \committeemember, per
% "Number of lines corresponds to number of committee members".
\renewcommand{\thesistitlepage}{%
	\loadgeometry{nopagenumber}%
	\begin{titlepage}%
		\thispagestyle{empty}\doublespacing\centering
		\thesis@tagtitle{Title}{\MakeUppercase{\@title}}\par
		By\par
		\MakeUppercase{\@author}\par
		A \thesis@docnoun{} submitted to the\par
		School of Graduate Studies\par
		Rutgers, The State University of New Jersey\par
		In partial fulfillment of the requirements\par
		For the degree of\par
		\thesis@degree\par
		Graduate Program in \thesis@program\par
		Written under the direction of\par
		\thesis@advisor\par
		And approved by\par
		\begingroup
		\renewcommand{\thesis@memberline}[2]{%
			\rule{3in}{0.4pt}\par
		}%
		\thesis@committee
		\endgroup
		New Brunswick, New Jersey\par
		\thesis@date\par
	\end{titlepage}%
	\loadgeometry{normal}%
}

% --- Committee approval page --------------------------------------------------
% Rutgers's filed document contains no separate approval page: approval is on
% the title page, as the blank signature lines above.
%
% The style guide fixes the order of sections and the list has no approval page
% in it -- "Starting with the Copyright page (which is optional) all the way to
% the Bibliography (a required section) the following sections are the order
% which must be followed in every dissertation/thesis" -- and the entries are
% Copyright, Title, Abstract, Acknowledgement/Dedication, Table of contents,
% List of tables, List of figures, List of illustrations.
%
% SIGNATURES ARE STILL COLLECTED ON PAPER, on a copy of the title page that is
% NOT the filed one: "The signed copy of this page gets submitted along with the
% rest of your checklist to sgs.degree.submissions@grad.rutgers.edu for review
% as per the checklist instructions. The title page within your entire
% dissertation which you will upload to ProQuest for publishing must not be
% signed."
\thesisnoapprovalpage
